% ===== PRÉAMBULE CANONIQUE — à coller VERBATIM en tête de CHAQUE .tex =====
\documentclass[11pt,a4paper]{article}
\usepackage[utf8]{inputenc}\usepackage[T1]{fontenc}\usepackage{lmodern}
\usepackage[french]{babel}
\usepackage{amsmath,amssymb,mathtools}\usepackage{icomma}
\usepackage[version=4]{mhchem}          % \ce{...} : équations & formules
\usepackage{siunitx}\sisetup{locale=FR} % \qty{}{}, \si{}, \ang{}
\usepackage{chemfig}                    % \chemfig{...} : structures développées
\usepackage{xcolor}\usepackage{colortbl}
\usepackage{tikz}\usepackage{pgfplots}\pgfplotsset{compat=1.18}
\usepackage[most]{tcolorbox}\usepackage{enumitem}\usepackage{array,booktabs}
\usepackage[a4paper,margin=2.2cm]{geometry}
\usetikzlibrary{arrows.meta,calc,angles,quotes,positioning,patterns,backgrounds,shapes.geometric}
\definecolor{primary}{RGB}{222,49,99}\definecolor{primarydark}{RGB}{150,22,55}
\definecolor{defcol}{RGB}{0,114,178}\definecolor{methcol}{RGB}{52,92,148}
\definecolor{excol}{RGB}{0,135,90}\definecolor{attncol}{RGB}{200,40,55}
\tcbset{boxbase/.style={enhanced,breakable,boxrule=0.9pt,arc=2.5pt,
  left=3mm,right=3mm,top=2.3mm,bottom=2.3mm,fonttitle=\bfseries,coltitle=white}}
% --- boîtes TUTORIEL (noms EXACTS, ne pas renommer) ---
\newtcolorbox{cle}[1][]{boxbase,colback=defcol!6,colframe=defcol,
  title={\textsc{À retenir}\ifx\relax#1\relax\relax\else\textmd{ : \itshape #1}\fi}}
\newtcolorbox{methode}[1][]{boxbase,colback=methcol!7,colframe=methcol,sharp corners,
  title={\textsc{Méthode}\ifx\relax#1\relax\relax\else\textmd{ --- \itshape #1}\fi}}
\newtcolorbox{exemplebox}[1][]{boxbase,colback=excol!5,colframe=excol,
  title={\textsc{Exemple}\ifx\relax#1\relax\relax\else\textmd{ : \itshape #1}\fi}}
\newtcolorbox{piege}[1][]{boxbase,colback=attncol!6,colframe=attncol,
  title={\textsc{Piège}\ifx\relax#1\relax\relax\else\textmd{ : \itshape #1}\fi}}
\newtcolorbox{remarque}[1][]{boxbase,colback=black!4,colframe=black!45,
  title={\textsc{Remarque}\ifx\relax#1\relax\relax\else\textmd{ : \itshape #1}\fi}}
% --- boîtes EXERCICE / CORRIGÉ (noms EXACTS, auto-comptées) ---
\newtcolorbox[auto counter]{exo}[1][]{boxbase,colback=primary!4,colframe=primary,
  title={\textsc{Exercice}~\thetcbcounter\ifx\relax#1\relax\relax\else\textmd{ --- \itshape #1}\fi}}
\newtcolorbox[auto counter]{corrige}[1][]{boxbase,colback=excol!4,colframe=excol,
  title={\textsc{Corrigé}~\thetcbcounter\ifx\relax#1\relax\relax\else\textmd{ --- \itshape #1}\fi}}
\newcommand{\R}{\mathbb{R}}\newcommand{\N}{\mathbb{N}}\newcommand{\Z}{\mathbb{Z}}\newcommand{\ds}{\displaystyle}
% =========================================================================
\begin{document}
\sisetup{per-mode=symbol}
\begin{exo}[écrire l'équation de dissolution]
Écris l'équation de dissolution dans l'eau de chacun de ces sels (avec les états $(s)$ et $(aq)$ et les
bons coefficients).
\begin{enumerate}[label=\alph*)]
  \item \ce{AgCl}
  \item \ce{CaCO3}
  \item \ce{CaF2}
  \item \ce{Fe(OH)3}
  \item \ce{Ag3PO4}
\end{enumerate}
\end{exo}

\begin{exo}[reconnaître le type $p{:}q$]
Pour chaque sel, donne le nombre de cations et d'anions libérés par une unité dissoute, puis son
type $p{:}q$ (cation$\,$:$\,$anion).
\begin{enumerate}[label=\alph*)]
  \item \ce{KCl}
  \item \ce{PbBr2}
  \item \ce{Ag2CrO4}
  \item \ce{Al(OH)3}
\end{enumerate}
\end{exo}

\begin{exo}[concentration des ions à saturation]
On note $s$ la solubilité molaire. Exprime $[\ce{cation}]$ et $[\ce{anion}]$ en fonction de $s$ à
saturation.
\begin{enumerate}[label=\alph*)]
  \item \ce{AgCl}
  \item \ce{CaF2}
  \item \ce{Ag3PO4}
\end{enumerate}
\end{exo}

\begin{exo}[du molaire au massique : $s_m = s\times M$]
Calcule la solubilité massique $s_m$ à partir de la solubilité molaire $s$ et de la masse molaire $M$.
\begin{enumerate}[label=\alph*)]
  \item \ce{AgCl} : $s=\qty{1.3e-5}{\mole\per\litre}$, $M=\qty{143.3}{\gram\per\mole}$.
  \item \ce{CaF2} : $s=\qty{2.1e-4}{\mole\per\litre}$, $M=\qty{78.1}{\gram\per\mole}$.
  \item \ce{BaSO4} : $s=\qty{1.0e-5}{\mole\per\litre}$, $M=\qty{233.4}{\gram\per\mole}$.
\end{enumerate}
\end{exo}

\begin{exo}[du massique au molaire : $s = s_m/M$]
Calcule la solubilité molaire $s$ à partir de la solubilité massique $s_m$ et de la masse molaire $M$.
\begin{enumerate}[label=\alph*)]
  \item \ce{CaCO3} : $s_m=\qty{6.9e-3}{\gram\per\litre}$, $M=\qty{100.1}{\gram\per\mole}$.
  \item \ce{PbBr2} : $s_m=\qty{4.84}{\gram\per\litre}$, $M=\qty{367.0}{\gram\per\mole}$.
\end{enumerate}
\end{exo}
\end{document}
